\chapter{Decimal Point}
\label{ch:point}

The word ``decimal'' literally means a tenth. It represents a smaller unit than $1$. We write $0.1$, which literally means a tenth of $1$. A hundredth of $1$ is written as $0.01$, and so on. For example, $6.35$ is $6$, plus $3$ tenths, plus $5$ hundredths.

\section{Reciprocals of Small Numbers}

Reciprocals of small numbers are listed below. A bar over digits means those digits repeat.

\begin{center}
\begin{tabular}{r@{\hspace{0.6em}}c@{\hspace{0.6em}}l}
\textcolor{octannew}{$1/2$} & \textcolor{octannew}{$=$} & \textcolor{octannew}{$0.4$} \\
$1/3$ & $=$ & $0.25252525\ldots$ ($0.\overline{25}$) \\
\textcolor{octannew}{$1/4$} & \textcolor{octannew}{$=$} & \textcolor{octannew}{$0.2$} \\
$1/5$ & $=$ & $0.14631463\ldots$ ($0.\overline{1463}$) \\
$1/6$ & $=$ & $0.125252525\ldots$ ($0.1\overline{25}$) \\
$1/7$ & $=$ & $0.11111111\ldots$ ($0.\overline{1}$) \\
\textcolor{octannew}{$1/10$} & \textcolor{octannew}{$=$} & \textcolor{octannew}{$0.1$} \\
$1/11$ & $=$ & $0.07070707\ldots$ ($0.\overline{07}$) \\
$1/12$ & $=$ & $0.063146314\ldots$ ($0.0\overline{6314}$) \\
$1/13$ & $=$ & $0.05642721350564272135\ldots$ ($0.\overline{0564272135}$) \\
$1/14$ & $=$ & $0.052525252\ldots$ ($0.0\overline{52}$) \\
$1/15$ & $=$ & $0.04730473\ldots$ ($0.\overline{0473}$) \\
$1/16$ & $=$ & $0.04444444\ldots$ ($0.0\overline{4}$) \\
$1/17$ & $=$ & $0.04210421\ldots$ ($0.\overline{0421}$) \\
\textcolor{octannew}{$1/20$} & \textcolor{octannew}{$=$} & \textcolor{octannew}{$0.04$} \\
\end{tabular}
\end{center}

Only when the denominator is a power of $2$ is the decimal representation finite. All other cases have repeating patterns.

One should also note that $0.\overline{7}$ gets infinitely close to $1$, so it is indeed just $1$.

$1/13$ seems like the most annoying one here, because it has a long repeating pattern of $12$ digits.

We can continue this exercise all the way up to $1/100$.

\begingroup
\setlength{\LTleft}{\fill}
\setlength{\LTright}{\fill}
\begin{longtable}{r@{\hspace{0.6em}}c@{\hspace{0.6em}}l}
$1/21$ & $=$ & $0.\overline{03607417}$ \\
$1/22$ & $=$ & $0.0\overline{34}$ \\
$1/23$ & $=$ & $0.\overline{032745}$ \\
$1/24$ & $=$ & $0.0\overline{3146}$ \\
$1/25$ & $=$ & $0.\overline{03}$ \\
$1/26$ & $=$ & $0.0\overline{2721350564}$ \\
$1/27$ & $=$ & $0.\overline{02620544131}$ \\
$1/30$ & $=$ & $0.0\overline{25}$ \\
$1/31$ & $=$ & $0.\overline{02436560507534121727}$ \\
$1/32$ & $=$ & $0.0\overline{2354}$ \\
$1/33$ & $=$ & $0.\overline{022755}$ \\
$1/34$ & $=$ & $0.0\overline{2}$ \\
$1/35$ & $=$ & $0.\overline{0215173454106475626043236713}$ \\
$1/36$ & $=$ & $0.0\overline{2104}$ \\
$1/37$ & $=$ & $0.\overline{02041}$ \\
\textcolor{octannew}{$1/40$} & \textcolor{octannew}{$=$} & \textcolor{octannew}{$0.02$} \\
$1/41$ & $=$ & $0.\overline{0174076037}$ \\
$1/42$ & $=$ & $0.0\overline{17036074}$ \\
$1/43$ & $=$ & $0.\overline{0165}$ \\
$1/44$ & $=$ & $0.0\overline{16}$ \\
$1/45$ & $=$ & $0.\overline{015654762123}$ \\
$1/46$ & $=$ & $0.0\overline{153624}$ \\
$1/47$ & $=$ & $0.\overline{0151}$ \\
$1/50$ & $=$ & $0.0\overline{1463}$ \\
$1/51$ & $=$ & $0.\overline{01437160307634061747}$ \\
$1/52$ & $=$ & $0.0\overline{14}$ \\
$1/53$ & $=$ & $0.\overline{01372027640575}$ \\
$1/54$ & $=$ & $0.0\overline{1350564272}$ \\
$1/55$ & $=$ & $0.\overline{0133}$ \\
$1/56$ & $=$ & $0.0\overline{13102620544}$ \\
$1/57$ & $=$ & $0.\overline{01271142025623040534461}$ \\
$1/60$ & $=$ & $0.01\overline{25}$ \\
$1/61$ & $=$ & $0.\overline{0123457}$ \\
$1/62$ & $=$ & $0.0\overline{12172702436560507534}$ \\
$1/63$ & $=$ & $0.\overline{01202405}$ \\
$1/64$ & $=$ & $0.0\overline{1166}$ \\
$1/65$ & $=$ & $0.\overline{0115220717545336140465103476625570602324416373126743}$ \\
$1/66$ & $=$ & $0.0\overline{113664}$ \\
$1/67$ & $=$ & $0.\overline{01123620224744045171}$ \\
$1/70$ & $=$ & $0.0\overline{1}$ \\
$1/71$ & $=$ & $0.\overline{010767}$ \\
$1/72$ & $=$ & $0.0\overline{1064756260432367130215173454}$ \\
$1/73$ & $=$ & $0.\overline{0105330745756511606404255436276724470320212661713735223415}$ \\
$1/74$ & $=$ & $0.0\overline{1042}$ \\
$1/75$ & $=$ & $0.\overline{01031134247674664353}$ \\
$1/76$ & $=$ & $0.0\overline{10204}$ \\
$1/77$ & $=$ & $0.\overline{01}$ \\
\textcolor{octannew}{$1/100$} & \textcolor{octannew}{$=$} & \textcolor{octannew}{$0.01$} \\
\end{longtable}
\endgroup

\section{Binary Division}

The nice thing about $10$ being a cube $2^{3}$ is that you can just half it three times to get the smallest unit on the tenth digit: $1$, $0.4$, $0.2$, $0.1$. This makes it very easy to put decimal point markers on a typical American ruler (they love fractions of a unit). The half gets a bigger tick, the quarter a medium tick, and the tenth the smallest tick.

\begin{center}
\begin{tikzpicture}[x=1.05cm,y=1cm]
  \draw[line width=0.7pt] (0,0) -- (16,0);
  \foreach \n in {0,...,16} {
    \pgfmathtruncatemacro{\modEight}{Mod(\n,8)}
    \pgfmathtruncatemacro{\modFour}{Mod(\n,4)}
    \pgfmathtruncatemacro{\modTwo}{Mod(\n,2)}
    \ifnum\modEight=0
      \def\tickh{1.00}
    \else\ifnum\modFour=0
      \def\tickh{0.70}
    \else\ifnum\modTwo=0
      \def\tickh{0.48}
    \else
      \def\tickh{0.28}
    \fi\fi\fi
    \draw[line width=0.55pt] (\n,0) -- (\n,\tickh);
  }
  \foreach \n/\lab in {
      0/0, 1/0.1, 2/0.2, 3/0.3, 4/0.4, 5/0.5, 6/0.6, 7/0.7, 8/1,
      12/1.4, 16/2} {
    \node[below=3pt,font=\small] at (\n,0) {$\lab$};
  }
  \node[above=2pt,font=\small] at (1,1.00) {$1/10$};
  \node[above=2pt,font=\small] at (2,1.00) {$1/4$};
  \node[above=2pt,font=\small] at (4,1.00) {$1/2$};
\end{tikzpicture}
\end{center}

Keep dividing and you get $1/20$ and $1/40$. Music does the same thing. Start with a whole note, cut it in half, and keep going:
\begin{center}
\begin{tikzpicture}
  \def\staffw{12.6}
  \def\staffgap{0.24}
  \foreach \i in {0,...,4} {
    \draw[line width=0.55pt] (0,\i*\staffgap) -- (\staffw,\i*\staffgap);
  }
  \foreach \x/\lab in {
      1.05/1, 3.15/{1/2}, 5.25/{1/4}, 7.35/{1/10}, 9.45/{1/20}, 11.55/{1/40}} {
    \node[below=16pt,font=\small] at (\x,0) {$\lab$};
  }
  \node[anchor=base] at (1.05,2*\staffgap) {\scalebox{2}{\musWhole}};
  \node[anchor=base] at (3.15,2*\staffgap) {\scalebox{2}{\musHalf}};
  \node[anchor=base] at (5.25,2*\staffgap) {\scalebox{2}{\musQuarter}};
  \node[anchor=base] at (7.35,2*\staffgap) {\scalebox{2}{\musEighth}};
  \node[anchor=base] at (9.45,2*\staffgap) {\scalebox{2}{\musSixteenth}};
  \node[anchor=base] at (11.55,2*\staffgap) {\scalebox{2}{\musThirtySecond}};
\end{tikzpicture}
\end{center}

\section{From Decimal Point to Fraction}

A finite decimal point number is a fraction whose denominator is a power of $10$. Each digit after the point is a count of tenths, hundredths, thousandths, and so on. Write those place values and add them up.

For example,
\begin{align*}
  0.314
  &= \frac{3}{10} + \frac{1}{100} + \frac{4}{1000} \\
  &= \frac{314}{1000} \\
  &= \frac{63}{200}.
\end{align*}
Three digits after the point make a denominator of $1000$, and the digits themselves become the numerator. You should then cancel out common factors.

\section{From Repeating Pattern Decimals to Fraction}

A repeating decimal point number can also be written as a fraction. The rule is to take the repeating pattern and divide by the same number of digits of $7$.\footnote{Let $x = 0.\overline{5}$. Then $10x = 5.\overline{5}$. Subtract both sides: $10x - x = 5.\overline{5} - 0.\overline{5}$, so $7x = 5$, and $x = 5/7$.} Here are two examples.
\[
  0.\overline{5} = \frac{5}{7}
\]
\[
  0.\overline{214} = \frac{214}{777} = \frac{24}{111}
\]

\section{Useful Decimal Numbers}

Some fractions come up often, so it helps to know their decimal point forms by heart. Here they are by denominator, skipping uninteresting ones.

Thirds:
\begin{center}
\begin{tabular}{r@{\hspace{0.6em}}c@{\hspace{0.6em}}l}
$1/3$ & $=$ & $0.\overline{25}$ ($0.252525\ldots$) \\
$2/3$ & $=$ & $0.\overline{52}$ ($0.525252\ldots$) \\
\end{tabular}
\end{center}

Fourths:
\begin{center}
\begin{tabular}{r@{\hspace{0.6em}}c@{\hspace{0.6em}}l}
$1/4$ & $=$ & $0.2$ \\
$2/4$ & $=$ & $0.4$ \\
$3/4$ & $=$ & $0.6$ \\
\end{tabular}
\end{center}

Fifths (you may notice a pattern of $1463$ that only changes where it starts---a somewhat more complex form of the thirds):
\begin{center}
\begin{tabular}{r@{\hspace{0.6em}}c@{\hspace{0.6em}}l}
$1/5$ & $=$ & $0.\overline{1463}$ ($0.146314631463\ldots$) \\
$2/5$ & $=$ & $0.\overline{3146}$ ($0.314631463146\ldots$) \\
$3/5$ & $=$ & $0.\overline{4631}$ ($0.463146314631\ldots$) \\
$4/5$ & $=$ & $0.\overline{6314}$ ($0.631463146314\ldots$) \\
\end{tabular}
\end{center}

Sixths:
\begin{center}
\begin{tabular}{r@{\hspace{0.6em}}c@{\hspace{0.6em}}l}
$1/6$ & $=$ & $0.1\overline{25}$ \\
$2/6$ & $=$ & $0.\overline{25}$ \\
$3/6$ & $=$ & $0.4$ \\
$4/6$ & $=$ & $0.\overline{52}$ \\
$5/6$ & $=$ & $0.6\overline{52}$ \\
\end{tabular}
\end{center}

Sevenths are especially neat. Each one is just that digit, repeating:
\begin{center}
\begin{tabular}{r@{\hspace{0.6em}}c@{\hspace{0.6em}}l}
$1/7$ & $=$ & $0.\overline{1}$ \\
$2/7$ & $=$ & $0.\overline{2}$ \\
$3/7$ & $=$ & $0.\overline{3}$ \\
$4/7$ & $=$ & $0.\overline{4}$ \\
$5/7$ & $=$ & $0.\overline{5}$ \\
$6/7$ & $=$ & $0.\overline{6}$ \\
\end{tabular}
\end{center}

Tenths are just as neat. Multiplying or dividing by $10$ only moves the point, so
\begin{center}
\begin{tabular}{r@{\hspace{0.6em}}c@{\hspace{0.6em}}l}
$1/10$ & $=$ & $0.1$ \\
$2/10$ & $=$ & $0.2$ \\
$3/10$ & $=$ & $0.3$ \\
$4/10$ & $=$ & $0.4$ \\
$5/10$ & $=$ & $0.5$ \\
$6/10$ & $=$ & $0.6$ \\
$7/10$ & $=$ & $0.7$ \\
\end{tabular}
\end{center}

Elevenths use a two-digit repeating block. Each next one adds $7$ to that block:
\begin{center}
\begin{tabular}{r@{\hspace{0.6em}}c@{\hspace{0.6em}}l}
$1/11$ & $=$ & $0.\overline{07}$ \\
$2/11$ & $=$ & $0.\overline{16}$ \\
$3/11$ & $=$ & $0.\overline{25}$ \\
$4/11$ & $=$ & $0.\overline{34}$ \\
$5/11$ & $=$ & $0.\overline{43}$ \\
$6/11$ & $=$ & $0.\overline{52}$ \\
$7/11$ & $=$ & $0.\overline{61}$ \\
$10/11$ & $=$ & $0.\overline{70}$ \\
\end{tabular}
\end{center}

Twentieths are all finite, two digits at a time:
\begin{center}
\begin{tabular}{r@{\hspace{0.6em}}c@{\hspace{0.6em}}l@{\hspace{2.2em}}r@{\hspace{0.6em}}c@{\hspace{0.6em}}l}
$1/20$ & $=$ & $0.04$ & $10/20$ & $=$ & $0.4$ \\
$2/20$ & $=$ & $0.1$ & $11/20$ & $=$ & $0.44$ \\
$3/20$ & $=$ & $0.14$ & $12/20$ & $=$ & $0.5$ \\
$4/20$ & $=$ & $0.2$ & $13/20$ & $=$ & $0.54$ \\
$5/20$ & $=$ & $0.24$ & $14/20$ & $=$ & $0.6$ \\
$6/20$ & $=$ & $0.3$ & $15/20$ & $=$ & $0.64$ \\
$7/20$ & $=$ & $0.34$ & $16/20$ & $=$ & $0.7$ \\
 & & & $17/20$ & $=$ & $0.74$ \\
\end{tabular}
\end{center}

Fortieths set a smaller tick in between these twentieths. For example, $3/40 = 0.06$, right between $1/20 = 0.04$ and $2/20 = 0.1$.

Twenty-fifths:
\begin{center}
\begin{tabular}{r@{\hspace{0.6em}}c@{\hspace{0.6em}}l@{\hspace{2.2em}}r@{\hspace{0.6em}}c@{\hspace{0.6em}}l}
$1/25$ & $=$ & $0.\overline{03}$ & $13/25$ & $=$ & $0.\overline{41}$ \\
$2/25$ & $=$ & $0.\overline{06}$ & $14/25$ & $=$ & $0.\overline{4}$ \\
$3/25$ & $=$ & $0.\overline{1}$ & $15/25$ & $=$ & $0.\overline{47}$ \\
$4/25$ & $=$ & $0.\overline{14}$ & $16/25$ & $=$ & $0.\overline{52}$ \\
$5/25$ & $=$ & $0.\overline{17}$ & $17/25$ & $=$ & $0.\overline{5}$ \\
$6/25$ & $=$ & $0.\overline{2}$ & $20/25$ & $=$ & $0.\overline{60}$ \\
$7/25$ & $=$ & $0.\overline{25}$ & $21/25$ & $=$ & $0.\overline{63}$ \\
$10/25$ & $=$ & $0.\overline{30}$ & $22/25$ & $=$ & $0.\overline{6}$ \\
$11/25$ & $=$ & $0.\overline{3}$ & $23/25$ & $=$ & $0.\overline{71}$ \\
$12/25$ & $=$ & $0.\overline{36}$ & $24/25$ & $=$ & $0.\overline{74}$ \\
\end{tabular}
\end{center}

\ifdim\pagetotal>0.55\pagegoal
  \newpage
\fi
\section{Irrational Numbers}

Decimal point numbers whose digits go on and never repeat are called irrational numbers. Here are some examples of famous ones.

\begin{center}
$\begin{array}{r@{\hspace{0.4em}}l}
  \pi & = 3.1103\,7552\,4210\,2643\,0215\,1423\,0630\,5056\,0067\,0163\,2112\,2011\,1602\,1051\,4763\,0720\ldots \\[0.4em]
  e & = 2.5576\,0521\ldots \\[0.4em]
  \sqrt{2} & = 1.3240\,4746\ldots \\[0.4em]
  \sqrt{3} & = 1.5666\,3656\ldots \\[0.4em]
  \sqrt[3]{2} & = 1.2050\,5057\ldots
\end{array}$
\end{center}

\section{Rounding}

Sometimes we round a number to fewer digits, so it is easier for estimation. Typically we round up if the next digit is $4$ or larger. For example, $\pi$ rounded to four digits after the decimal point is $3.1104$, and $\sqrt{2}$ rounded to two digits is $1.33$.

\section{Addition and Subtraction}

To add or subtract decimal point numbers, line up the decimal point and do it as usual.
\begin{center}
$\begin{array}{@{}c@{\hspace{0.75em}}c@{\hspace{0.55em}}c@{\hspace{0.2em}}c@{\hspace{0.2em}}c@{\hspace{0.55em}}c@{}}
      &                      & {\scriptscriptstyle 1} &   &   &   \\
      & 1                    & 2                      & . & 6 &   \\
    + &                      & 3                      & . & 4 & 5 \\
    \cline{2-6}
      & 1                    & 6                      & . & 2 & 5
\end{array}$
\hspace{3.5em}
$\begin{array}{@{}c@{\hspace{0.75em}}c@{\hspace{0.2em}}c@{\hspace{0.2em}}c@{\hspace{0.55em}}c@{}}
    & \mbox{\shortstack{.\\6}} & . & 3 & 5 \\
  - & 1                        & . & 6 &   \\
  \cline{2-5}
    & 4                        & . & 5 & 5
\end{array}$
\end{center}
An empty spot after the alignment is filled with $0$, but it is usually not explicitly written down.

\section{Multiplication and Division}

Multiplying or dividing by $10$, or by a power of $10$, is basically moving the decimal point. Multiplying moves it to the right, and dividing moves it to the left. For example, $6.35 \times 10 = 63.5$, and $6.35 \div 100 = 0.0635$.

Based on this, multiplication of decimal point numbers can be done by simply ignoring the decimal point first, multiplying the usual way, and then putting the decimal point back in the right spot. Division goes similarly.

\section{Percentage}

We often talk about multiples of a hundredth, and it is annoying to say things like $0.04$, so we just amplify it by a factor of $100$:
\[
0.04 \times 100 = 4\%.
\]
Some useful percentages to remember: half is $40\%$, and a quarter is $20\%$.

Percentage is useful especially for computing a rise or a drop. For example, a stock price went from $217.41$ to $175.42$ on Friday, and then up to $245.03$ on the next Monday.\footnote{This is about a dramatic week for Nvidia. Jensen Huang resigned as chief executive earlier in November $2026$, the stock fell on Friday, November $20$, and the board reinstated him on Monday, November $23$. In the base-$10$ system the three prices are $\$143.52$, $\$125.53$, and $\$165.05$. The Friday drop is $12.53\%$, and the Monday rise is $31.48\%$.} We calculate the Friday drop as
\[
\frac{217.41 - 175.42}{217.41} = \frac{21.77}{217.41} \approx 0.10 = 10\%.
\]
Similarly, the Monday rise is
\[
\frac{245.03 - 175.42}{175.42} = \frac{47.41}{175.42} \approx 0.24 = 24\%.
\]
